\Needspace{29\baselineskip}
\subsection{Comparison of the reduced Planck constant}
\label{sec:revision-planck-contraste}

The International System in force since 20 May 2019 fixes
\(h=6.62607015\times10^{-34}\,\mathrm{J\,s}\) exactly
\cite{bipm2018}. Its reduced constant is therefore
\[
 \hbar_{\mathrm{SI}}=\frac{6.62607015}{2\pi}\,10^{-34}\,\mathrm{J\,s}
 =1.0545718176461563912\ldots\times10^{-34}\,\mathrm{J\,s}.
\]
The ellipsis denotes the expansion of an exact expression,
not an experimental uncertainty in \(h\) in the current SI.
This reference value is used here after evaluation of the
HMT readers; its function is comparative.

The article's two return coordinates are
\begin{align*}
 H_5&=1.0545718176461544696\ldots,\\
 \eta_{\mathrm{ret}}&=1.0545718169150390611\ldots.
\end{align*}
Under the identification of \(\mathcal U_S\) with \(\mathrm{J\,s}\), the
relative difference between the preceding section and
\(\hbar_{\mathrm{SI}}\) is approximately \(-1.82\times10^{-15}\).
For the subsequent section, it is approximately \(-6.93\times10^{-10}\).
The second difference contains the return effect used by the
electronic corrector; it does not equal the evaluation error of the
first section.

The precision of these comparisons allows the quantitative proximity
of the construction of \(H_5\) to be recognized while distinguishing
proximity from exact identity. In the present manuscript, the displayed
expressions do not provide an exact equality with \(h/(2\pi)\).
The result proved for \(\Sigma_H\) is its decade \(-34\);
the analytic result for \(H_5\) is the evaluation of the stated
character. The physical identification is examined in this chart and with the
explicit differences just given. The SI value does not participate
in the definition of the determinantal norm or in the regional recurrence.

